\scanpage{191}
où $E$ est le fibré vectoriel $\widetilde C$ sur $X$, et $\widehat E$ son complété formel le long de $X$, donc
\[
\operatorname{Pic}(\widehat E)\simeq\deleted{\unsure{\operatorname{Pic}}}\varprojlim\operatorname{Pic}(E_n)
\]
où $E_n$ est le voisinage infinitésimal d'ordre $n$ de $X$ dans $E$,
\[
E_n=(X,\underline O_E/J^{n+1})
\]
$J$ étant l'idéal définissant $X$ dans $E$. Notons ici
\[
J/J^2=\underline O_X(1)
\]
\[
\operatorname{Pic}(E)\xleftarrow{\sim}\operatorname{Pic}(X)
\]
\inserted{Donc $\operatorname{Pic}(S-a)\mathbin{-}\operatorname{Pic}(\widehat S-\hat a)$ est inj (surj) sss $\operatorname{Pic}(E)\mathbin{-}\operatorname{Pic}(\widehat E)$ l'est.}

% slanted
On sait déjà \inserted{au moins dans le cas \deleted{\unsure{où}} $\unsure{\widehat S|S}$ normal\ldots.} \deleted{\unsure{que}} (et on voit immédiatement sur la dernière formule) que $\operatorname{Pic}(S-a)\mathbin{-}\operatorname{Pic}(\widehat S-\hat a)$ est \uline{injectif} donc $\operatorname{Pic}(E)\longrightarrow\operatorname{Pic}(\widehat E)$ est injectif [car $\operatorname{Pic}(X)\mathbin{-}\operatorname{Pic}(\widehat E)$ l'est, ayant l'inverse à gauche $\operatorname{Pic}(\widehat E)\mathbin{-}\operatorname{Pic}(E_0)=\operatorname{Pic}(X)$]

Reste à voir si $\operatorname{Pic}(E)\mathbin{-}\operatorname{Pic}(\widehat E)$ est bijectif. Grâce à $\operatorname{Pic}(X)\unsure{\simeq|\xrightarrow{\sim}}\operatorname{Pic}(E)$
% archive sheet number: 190
